\section{线段树}

\subsection{静态线段树}

支持区间修改、区间查询。\texttt{Info} 与 \texttt{Tag} 需按题目自实现：\texttt{Info} 存区间聚合值，须带区间端点成员 \texttt{l,r}（树自动写入，供 \texttt{apply} 判断作用长度），重载 \texttt{operator+} 合并两个相邻子区间、提供 \texttt{apply(Tag)} 作整段修改；\texttt{Tag} 为懒标记，需提供 \texttt{apply(Tag)} 叠加标记，默认构造表示无待下传标记。

构造 \texttt{SegmentTree<Info,Tag>(n)} 得 $[1,n]$ 上的全零树，或传下标 $1..n$ 的初值数组（大小 $n+1$）建树，\texttt{rebuild} 可重新初始化。\texttt{modify(l,r,tag)} 对 $[l,r]$ 整段打标记，\texttt{seg(l,r)}（\texttt{operator()}）返回聚合值，越界下标自动夹到 $[1,n]$。单次 $\mathcal{O}(\log n)$。

\begin{lstlisting}
template<typename Info,typename Tag>
struct SegmentTree{
private:
    int _n = 0;
    std::vector<Info> _info;
    std::vector<Tag> _tag;
public:
    SegmentTree(int n = 0){
        rebuild(n);
    }
    SegmentTree(const std::vector<Info>& v){
        assert(v.size() >= 1);
        rebuild(v.size() - 1,v);
    }
    void rebuild(int n,const std::vector<Info>& v = std::vector<Info>(0)){
        _n = n;
        _info.assign((_n + 2) << 2,Info());
        _tag.assign((_n + 2) << 2,Tag());
        if (n > 0) build(1,1,n,v);
    }
private:
    void build(int i,int l,int r,const std::vector<Info>& v){
        if(l == r){
            if (v.size() > l) _info[i] = v[l];
            resize(_info[i],l,r);
            return;
        }
        int mid = l + r >> 1;
        build(i << 1,l,mid,v);
        build(i << 1 | 1,mid + 1,r,v);
        push_up(i,l,r);
    }
    void resize(Info& info,int l,int r){
        info.l = l;
        info.r = r;
    }
    void push_down(int i){
        _info[i << 1].apply(_tag[i]);
        _info[i << 1 | 1].apply(_tag[i]);
        _tag[i << 1].apply(_tag[i]);
        _tag[i << 1 | 1].apply(_tag[i]);
        _tag[i] = Tag();
    }
    void push_up(int i,int l,int r){
        _info[i] = _info[i << 1] + _info[i << 1 | 1];
        resize(_info[i],l,r);
    }
    void modify(int i,const Tag& val,int l,int r,int ql,int qr){
        if (ql <= l && r <= qr) {
            _info[i].apply(val);
            _tag[i].apply(val);
            return;
        }
        push_down(i);
        int mid = l + r >> 1;
        if (ql <= mid) modify(i << 1, val, l, mid, ql, qr);
        if (qr > mid) modify(i << 1 | 1, val, mid+1, r, ql, qr);
        push_up(i,l,r);
    }
    Info ask(int i,int l,int r,int ql,int qr){
        if (ql <= l && r <= qr) {
            return _info[i];
        }
        push_down(i);
        int mid = l + r >> 1;
        if (qr <= mid) return ask(i << 1, l, mid, ql, qr);
        else if (ql > mid)return ask(i << 1 | 1, mid+1, r, ql, qr);
        Info res = ask(i << 1, l, mid, ql, qr);
        res = res + ask(i << 1 | 1, mid+1, r, ql, qr);
        resize(res,std::max(ql,l),std::min(qr,r));
        return res;
    }
public:
    void modify(int l,int r,const Tag& val){
        if (l > r) return;
        modify(1,val,1,_n,std::max(1,l),std::min(r,_n));
    }
    Info operator()(int l,int r){
        if (l > r) return Info();
        return ask(1,1,_n,std::max(1,l),std::min(r,_n));
    }
};

struct Tag{
    Tag(){}
    void apply(const Tag& t){
        
    }
};

struct Info{
    int l = 1,r = 0;
    
    Info(){}
    void apply(const Tag& t){
        
    }
    
    friend Info operator+(const Info& lhs,const Info& rhs){
        Info res;
        return res;
    }
};
\end{lstlisting}

\subsection{开点线段树}

\noindent 值域 $[1,n]$ 极大（如 $10^{18}$）或整体稀疏时使用：只给被访问到的区间分配结点，单次修改/查询新建 $\mathcal{O}(\log n)$ 个结点。\texttt{Info}、\texttt{Tag} 的实现要求与静态版相同（此结构与静态版同名，同时使用需改名）。

构造 \texttt{SegmentTree<Info,Tag>(n)}（$n$ 为 \texttt{long long}），\texttt{modify(l,r,tag)} 对 $[l,r]$ 整段打标记，\texttt{seg(l,r)}（\texttt{operator()}）返回 $[l,r]$ 的聚合值；此处不做越界夹取，要求 $1 \le l \le r \le n$。单次 $\mathcal{O}(\log n)$。

\begin{lstlisting}
template<typename Info,typename Tag>
struct SegmentTree{
private:
    long long _n = 0;
    int _cur = 0,_root = 0;
    std::vector<Info> _info;
    std::vector<Tag> _tag;
    std::vector<int> _lc,_rc;

    void resize(Info& info,long long l,long long r){
        info.l = l;
        info.r = r;
    }
    int new_node(long long l,long long r){
        _cur++;
        _info.push_back(Info());
        _tag.push_back(Tag());
        _lc.push_back(-1);
        _rc.push_back(-1);
        resize(_info[_cur],l,r);
        return _cur;
    }
    void push_down(int u,long long l,long long r,long long mid) {
        if(_lc[u]<0){
            _lc[u] = new_node(l,mid);
        }
        if(_rc[u]<0){
            _rc[u] = new_node(mid + 1,r);
        }
        _info[_lc[u]].apply(_tag[u]);
        _tag[_lc[u]].apply(_tag[u]);
        _info[_rc[u]].apply(_tag[u]);
        _tag[_rc[u]].apply(_tag[u]);
        _tag[u] = Tag();
    }
    void push_up(int u,long long l,long long r){
        _info[u] = _info[_lc[u]] + _info[_rc[u]];
        resize(_info[u],l,r);
    }
    void modify(int u,const Tag& val,long long l,long long r,long long ql,long long qr){
        if (u<0) u = new_node(l,r);
        if (ql <= l && r <= qr) {
            _info[u].apply(val);
            _tag[u].apply(val);
            return;
        }
        long long mid = l + r >> 1;
        push_down(u,l,r,mid);
        if (ql <= mid) modify(_lc[u], val, l, mid, ql, qr);
        if (qr > mid) modify(_rc[u], val, mid+1, r,ql, qr);
        push_up(u,l,r);
    }
    Info ask(int u,long long l,long long r,long long ql,long long qr){
        if (u<0) u = new_node(l,r);
        if (ql <= l && r <= qr) {
            return _info[u];
        }
        long long mid = l + r >> 1;
        push_down(u,l,r,mid);
        if (qr <= mid) return ask(_lc[u], l, mid, ql, qr);
        else if (ql > mid)return ask(_rc[u], mid+1, r, ql, qr);
        Info res = ask(_lc[u], l, mid, ql, qr);
        res = res + ask(_rc[u], mid+1, r, ql, qr);
        resize(res,std::max(ql,l),std::min(qr,r));
        return res;
    }
public:
    SegmentTree(long long n){
        _cur = _root = -1;
        _n = n;
        _root = new_node(1,_n);
    }
    void modify(long long l,long long r,const Tag& val){
        if (l > r) return;
        modify(_root,val,1,_n,l,r);
    }
    Info operator()(long long l,long long r){
        if (l > r) return Info();
        return ask(_root,1,_n,l,r);
    }
};

struct Tag{
    
    Tag(){}
    void apply(const Tag& t){
        
    }
};

struct Info{
    long long l = 1,r = 0;
    
    Info(){}
    
    void apply(const Tag& t){
        
    }
    
    friend Info operator+(const Info& lhs,const Info& rhs){
        Info res;
        
        return res;
    }
};
\end{lstlisting}

\subsection{可持久化线段树}

支持保留历史版本，各版本可独立查询、继续修改。\texttt{Info}、\texttt{Tag} 的实现要求与静态版相同。

构造 \texttt{PersistentSegmentTree<Info,Tag>(n)} 或调用 \texttt{rebuild(n)} 建立 $0$ 号版本（空树）。\texttt{modify(ver,l,r,tag)} 以第 \texttt{ver} 个版本为基础修改并返回新版本编号（旧版本不变），之后用返回的编号继续操作即形成版本链；\texttt{seg(ver,l,r)}（\texttt{operator()}）查询第 \texttt{ver} 个版本的 $[l,r]$，返回 \texttt{Info}。不做标记永久化，查询时沿路径累计标记，故单次时空均为 $\mathcal{O}(\log n)$，$q$ 次操作总计 $\mathcal{O}((n+q)\log n)$。

\begin{lstlisting}
template<typename Info,typename Tag>
struct PersistentSegmentTree {
    int _n = 0;
    int _cur = 0;
    std::vector<int> roots;
    std::vector<Info> _info;
    std::vector<Tag> _tag;
    std::vector<int> _lc;
    std::vector<int> _rc;
    void resize(Info& info,int l,int r){
        info.l = l;
        info.r = r;
    }
    int new_node(int u,int l=-1,int r=-1){
        int v = _cur++;
        _info.emplace_back(_info[u]);
        _tag.emplace_back(_tag[u]);
        _lc.push_back(u>0?_lc[u]:0);
        _rc.push_back(u>0?_rc[u]:0);
        if(u==0)resize(_info[v],l,r);
        return v;
    }
    void push_down(int u,int l,int r,int mid) {
        _lc[u] = new_node(_lc[u],l,mid);
        _rc[u] = new_node(_rc[u],mid+1,r);
        _info[_lc[u]].apply(_tag[u]);
        _tag[_lc[u]].apply(_tag[u]);
        _info[_rc[u]].apply(_tag[u]);
        _tag[_rc[u]].apply(_tag[u]);
        _tag[u] = Tag();
    }
    void push_up(int u){
        _info[u] = _info[_lc[u]] + _info[_rc[u]];
    }
    void modify(int u,const Tag& val,int l,int r,int ql,int qr){
        int v = u;
        if (ql <= l && r <= qr) {
            _info[v].apply(val);
            _tag[v].apply(val);
            return;
        }
        int mid = l + r >> 1;
        push_down(v,l,r,mid);
        if (ql <= mid)modify(_lc[v],val,l,mid,ql,qr);
        if (qr  > mid)modify(_rc[v],val,mid+1,r,ql,qr);
        push_up(v);
        resize(_info[v],l,r);
        return;
    }
    Info ask(int u,Tag val,int l,int r,int ql,int qr){
        if (u == 0) {
            Info res;
            resize(res,std::max(ql,l),std::min(qr,r));
            res.apply(val);
            return res;
        }
        if (ql <= l && r <= qr) {
            Info res = _info[u];
            res.apply(val);
            return res;
        }
        val.apply(_tag[u]);
        int mid = l + r >> 1;
        if (qr <= mid) {
            return ask(_lc[u],val,l,mid,ql,qr);
        } else if (ql > mid) {
            return ask(_rc[u],val,mid+1,r,ql,qr);
        } else {
            Info left = ask(_lc[u],val,l,mid,ql,qr);
            Info right = ask(_rc[u],val,mid+1,r,ql,qr);
            Info res = left + right;
            resize(res,std::max(ql,l),std::min(qr,r));
            return res;
        }
    }
public:
    PersistentSegmentTree(int n = 0){
        rebuild(n);
    }
    void rebuild(int n){
        _n = n;
        _cur = 1;
        roots.clear();
        _info.assign(1,Info());
        _tag.assign(1,Tag());
        _lc.assign(1,0);
        _rc.assign(1,0);
        roots.push_back(new_node(0,1,_n));
    }
    int modify(int ver,int l,int r,const Tag& val){
        if (l > r) return ver;
        int new_root = new_node(roots[ver]);
        modify(new_root,val,1,_n,l,r);
        roots.push_back(new_root);
        return (int)roots.size()-1;
    }
    Info operator()(int ver,int l,int r){
        if (l > r) return Info();
        return ask(roots[ver],Tag(),1,_n,l,r);
    }
};
struct Tag{
    
    Tag(){}
    void apply(const Tag& t){

    }
};
struct Info{
    int l = 1,r = 0;

    Info(){}
    void apply(const Tag& t){
    }
    friend Info operator+(const Info& lhs,const Info& rhs){
        Info res;
        return res;
    }
};
\end{lstlisting}

\subsection{线段树合并与分裂}

动态开点线段树，初始为空树，支持把两棵树的同位置元素叠加合并、以及按位置区间把一棵树拆成两棵，常用于「每个点/每种颜色一棵权值线段树」之后自底向上合并。\texttt{Info}、\texttt{Tag} 的实现要求与静态版相同；两棵树各自的未下传标记只作用于自己的元素，合并时不会串到对方上。

构造 \texttt{MergeSplitSegmentTree<Info,Tag>(n)}，值域为 $[1,n]$（\texttt{long long}）。根编号为 \texttt{int}，$0$ 表示空树；\texttt{modify(root,p,val)} 把 \texttt{Info} 类型的 \texttt{val} 叠加到下标 $p$（插入元素用），\texttt{modify(root,l,r,tag)} 对区间打标记，两者都以引用传入 \texttt{root} 并就地改为新根；\texttt{root(l,r)}（\texttt{operator()}）返回 $[l,r]$ 的聚合值；\texttt{merge(a,b)} 把 $b$ 并入 $a$ 并返回新根（$b$ 不再可用，调用后应置 $0$），单次代价为两棵树的重合结点数，因此按任意顺序合并的总代价不超过历史新建的结点数；\texttt{split(root,l,r)} 把 $[l,r]$ 内的元素拆出作为新树返回，\texttt{root} 保留其余部分（区间内无元素时返回 $0$），单次 $\mathcal{O}(\log n)$。

\begin{lstlisting}
template<typename Info,typename Tag>
struct MergeSplitSegmentTree{
private:
    long long _n = 0;
    std::vector<Info> _info;
    std::vector<Tag> _tag;
    std::vector<int> _lc,_rc;
    void resize(Info& info,long long l,long long r){
        info.l = l;
        info.r = r;
    }
    int newNode(long long l,long long r){
        int u = _info.size();
        _info.push_back(Info());
        _tag.push_back(Tag());
        _lc.push_back(0);
        _rc.push_back(0);
        resize(_info[u],l,r);
        return u;
    }
    void push_down(int u,long long l,long long r,long long mid){
        if (_lc[u] == 0) _lc[u] = newNode(l,mid);
        if (_rc[u] == 0) _rc[u] = newNode(mid+1,r);
        _info[_lc[u]].apply(_tag[u]);
        _tag[_lc[u]].apply(_tag[u]);
        _info[_rc[u]].apply(_tag[u]);
        _tag[_rc[u]].apply(_tag[u]);
        _tag[u] = Tag();
    }
    void push_up(int u,long long l,long long r){
        _info[u] = _info[_lc[u]] + _info[_rc[u]];
        resize(_info[u],l,r);
    }
    void push_tag(int u,const Tag& val){
        if (u == 0) return;
        _info[u].apply(val);
        _tag[u].apply(val);
    }
    int modify(int u,long long l,long long r,long long p,const Info& v){
        if (u == 0) u = newNode(l,r);
        if (l == r) {
            _info[u] = _info[u] + v;
            resize(_info[u],l,r);
            return u;
        }
        long long mid = l + r >> 1;
        push_down(u,l,r,mid);
        if (p <= mid) _lc[u] = modify(_lc[u],l,mid,p,v);
        else _rc[u] = modify(_rc[u],mid+1,r,p,v);
        push_up(u,l,r);
        return u;
    }
    void modify(int u,const Tag& val,long long l,long long r,long long ql,long long qr){
        if (ql <= l && r <= qr) {
            _info[u].apply(val);
            _tag[u].apply(val);
            return;
        }
        long long mid = l + r >> 1;
        push_down(u,l,r,mid);
        if (ql <= mid) modify(_lc[u],val,l,mid,ql,qr);
        if (qr > mid) modify(_rc[u],val,mid+1,r,ql,qr);
        push_up(u,l,r);
    }
    Info ask(int u,Tag val,long long l,long long r,long long ql,long long qr){
        if (u == 0) {
            Info res;
            resize(res,std::max<long long>(ql,l),std::min<long long>(qr,r));
            res.apply(val);
            return res;
        }
        if (ql <= l && r <= qr) {
            Info res = _info[u];
            res.apply(val);
            return res;
        }
        val.apply(_tag[u]);
        long long mid = l + r >> 1;
        if (qr <= mid) return ask(_lc[u],val,l,mid,ql,qr);
        if (ql > mid) return ask(_rc[u],val,mid+1,r,ql,qr);
        Info res = ask(_lc[u],val,l,mid,ql,qr) + ask(_rc[u],val,mid+1,r,ql,qr);
        resize(res,std::max<long long>(ql,l),std::min<long long>(qr,r));
        return res;
    }
    int merge(int a,int b,long long l,long long r){
        if (a == 0) return b;
        if (b == 0) return a;
        if (l == r) {
            _info[a] = _info[a] + _info[b];
            _tag[a] = Tag();
            return a;
        }
        long long mid = l + r >> 1;
        push_down(a,l,r,mid);
        push_down(b,l,r,mid);
        _lc[a] = merge(_lc[a],_lc[b],l,mid);
        _rc[a] = merge(_rc[a],_rc[b],mid+1,r);
        push_up(a,l,r);
        return a;
    }
    int split(int u,long long l,long long r,long long ql,long long qr,int& rest){
        if (u == 0 || qr < l || r < ql) {
            rest = u;
            return 0;
        }
        if (ql <= l && r <= qr) {
            rest = 0;
            return u;
        }
        long long mid = l + r >> 1;
        int left = 0,right = 0;
        int a = split(_lc[u],l,mid,ql,qr,left);
        int b = split(_rc[u],mid+1,r,ql,qr,right);
        if (a == 0 && b == 0) {
            rest = u;
            return 0;
        }
        _lc[u] = left;
        _rc[u] = right;
        push_tag(left,_tag[u]);
        push_tag(right,_tag[u]);
        push_tag(a,_tag[u]);
        push_tag(b,_tag[u]);
        _tag[u] = Tag();
        push_up(u,l,r);
        int v = newNode(l,r);
        _lc[v] = a;
        _rc[v] = b;
        push_up(v,l,r);
        rest = u;
        return v;
    }
public:
    MergeSplitSegmentTree(long long n = 0){
        _n = n;
        _info.assign(1,Info());
        _tag.assign(1,Tag());
        _lc.assign(1,0);
        _rc.assign(1,0);
    }
    void modify(int& root,long long p,const Info& val){
        if (root == 0) root = newNode(1,_n);
        root = modify(root,1,_n,p,val);
    }
    void modify(int& root,long long l,long long r,const Tag& val){
        if (l > r) return;
        if (root == 0) root = newNode(1,_n);
        modify(root,val,1,_n,l,r);
    }
    Info operator()(int root,long long l,long long r){
        if (l > r) return Info();
        return ask(root,Tag(),1,_n,l,r);
    }
    int merge(int a,int b){
        return merge(a,b,1,_n);
    }
    int split(int& root,long long l,long long r){
        if (l > r) return 0;
        int rest = 0;
        int v = split(root,1,_n,l,r,rest);
        root = rest;
        return v;
    }
};
struct Tag{

    Tag(){}
    void apply(const Tag& t){

    }
};
struct Info{
    long long l = 1,r = 0;

    Info(){}
    void apply(const Tag& t){

    }
    friend Info operator+(const Info& lhs,const Info& rhs){
        Info res;
        return res;
    }
};
\end{lstlisting}
